\documentclass[11pt]{article}
\usepackage{amsmath,amssymb,amsthm,hyperref}
\newtheorem{theorem}{Theorem}
\title{An affine-relation obstruction to discretizing a second die}
\author{Methodological research note; independently reviewed}
\date{30 September 2026}
\begin{document}\maketitle
\begin{theorem}
Let $A$ be supported on finitely many distinct points
$a_1,\ldots,a_K\in[0,1]$, with positive rational probabilities
$w_1,\ldots,w_K$. Suppose $\sum_iw_i a_i=1/2$ and that
\[
 \left\{c\in\mathbb Q^K:\sum_i c_i a_i\in\mathbb Q\right\}
 =\mathbb Q\,(w_1,\ldots,w_K).
\]
There is no finitely supported random variable $B$ with rational
probabilities such that the independent variables $A,B,U$, with
$U$ uniform on $[0,1]$, have each of their six strict orders with
probability $1/6$.
\end{theorem}
\begin{proof}
Put $b_i=\Pr(B>a_i)\in\mathbb Q$. The order $U<A<B$ would give
\[
 \sum_iw_i a_i b_i=\frac16.
\]
The affine-relation hypothesis then forces
$(w_i b_i)_i$ to be a rational multiple of $(w_i)_i$.
Thus all $b_i$ have one common value $b$.
The six strict order probabilities sum to one, so $A$ and $B$ have
no ties. Marginalizing their relative order gives
$b=\Pr(A<B)=1/2$. But then
\[
 \Pr(U<A<B)=b\sum_iw_i a_i=\frac14,
\]
contradicting $1/6$.
\end{proof}

For equiprobable $A$, the hypothesis says that the only rational
coefficient combinations of its node values that are rational are
constant multiples of their sum. This note does not claim that all
real interval quadratures have this property, or that two small finite
dice are impossible. It identifies an arithmetic obstruction to a
construction strategy which first selects a real scalar quadrature and
then tries to discretize another variable while leaving the remaining
variables exactly uniform. Special affine relations among the first
quadrature's nodes are necessary for that strategy.
\end{document}
